% !TEX program = xelatex
\documentclass[11pt,a4paper]{article}

\usepackage{fontspec}
\usepackage{polyglossia}
\usepackage{amsmath,amssymb,amsfonts}
\usepackage{xcolor}
\usepackage{graphicx}
\usepackage{geometry}

\geometry{left=2cm,right=2cm,top=2cm,bottom=2cm}

\setdefaultlanguage{persian}
\setotherlanguage{english}
\setmainfont{XB Niloofar}

\newcommand{\en}[1]{\textenglish{#1}}
\newcommand{\separator}{\vspace{1em}\noindent\rule{\linewidth}{0.4pt}\vspace{1em}}
\newcommand{\thickseparator}{\vspace{1em}\noindent\rule{\linewidth}{1pt}\vspace{1em}}

\title{\textbf{پاسخ سوال ۳} \\[0.4cm] \small محاسبه گرادیان نسبت به بردارهای \en{Embedding}}
\author{محمد باغبان زاده \\ شماره دانشجویی: شناسهٔ دانشجو}
\date{}
\linespread{1.35}

\begin{document}

\maketitle
\thickseparator

\section{مفاهیم پایه: \en{Backpropagation} تا لایه \en{Embedding}}

بر اساس بخش‌های \textbf{\en{۶-۵}} و \textbf{\en{۶-۶}} کتاب، برای به‌روزرسانی بردارهای \en{Embedding} باید گرادیان تابع هزینه را نسبت به آن‌ها محاسبه کنیم. در این سوال، وزن‌های شبکه (\(W^{(1)}\) و \(W^{(2)}\)) ثابت فرض شده‌اند و فقط \en{Embedding}ها به‌روزرسانی می‌شوند.

\separator
\section{داده‌های مسأله (از سوال قبل)}

\subsection{ماتریس \en{Embedding} (E)}
\[
E = 
\begin{bmatrix}
0.2 & -0.1 \\
0.5 & 0.3 \\
-0.4 & 0.2 \\
0.1 & 0.0
\end{bmatrix}
\]

\subsection{وزن‌های شبکه}
\[
W^{(1)} = 
\begin{bmatrix}
0.8 & -0.2 \\
0.1 & 0.6
\end{bmatrix}
\quad,\quad
W^{(2)} = 
\begin{bmatrix}
0.5 \\
-0.3
\end{bmatrix}
\]

\subsection{ایندکس کلمات و \en{Embedding} سند (از سوال قبل)}
\begin{align*}
\text{:کلمه ۱ (ایندکس ۱)} &\quad e_1 = [0.5, 0.3] \\
\text{:کلمه ۲ (ایندکس ۲)} &\quad e_2 = [-0.4, 0.2] \\
\text{:\en{Embedding} سند (میانگین)} &\quad e_{\text{doc}} = [0.05, 0.25]
\end{align*}

\subsection{نتایج \en{Forward Pass} (از سوال قبل)}
\begin{align*}
h_{\text{net}} &= [-0.01, 0.155] \\
h &= [0, 0.155] \\
y_{\text{net}} &= -0.0465 \\
\hat{y} &\approx 0.4884
\end{align*}

\separator
\section{فرض برچسب واقعی}

مانند سوال قبل، فرض می‌کنیم برچسب واقعی \(y = 1\) است.

\separator
\section{گام ۱: محاسبه سیگنال خطا در لایه خروجی}

بر اساس بخش \textbf{\en{۶-۶-۴}} کتاب (معادله \en{۶-۴۳})، مشتق تابع هزینه نسبت به \(y_{\text{net}}\) برای تابع سیگموئید و کراس-آنتروپی به صورت زیر است:

\[
\delta^{(2)} = \frac{\partial L}{\partial y_{\text{net}}} = \hat{y} - y
\]

\[
\delta^{(2)} = 0.4884 - 1 = -0.5116
\]

\separator
\section{گام ۲: انتشار خطا به لایه پنهان}

سیگنال خطا برای هر نورون لایه پنهان به صورت زیر محاسبه می‌شود:

\[
\delta^{(1)}_j = \delta^{(2)} \cdot W^{(2)}_j \cdot \text{ReLU}'(h_{\text{net},j})
\]

که در آن \(\text{ReLU}'(x)\) مشتق تابع \en{ReLU} است (بخش \en{۶-۶-۴} کتاب، معادله \en{۶-۴۰}):

\[
\text{ReLU}'(x) = 
\begin{cases} 
1 if  x > 0 \\
0 if  x \le 0
\end{cases}
\]

\begin{align*}
\text{برای } j=1: &\quad h_{\text{net},1} = -0.01 \le 0 \quad \Rightarrow \quad \text{ReLU}'(h_{\text{net},1}) = 0 \\
&\quad \delta^{(1)}_1 = (-0.5116) \times 0.5 \times 0 = 0
\end{align*}

\begin{align*}
\text{برای } j=2: &\quad h_{\text{net},2} = 0.155 > 0 \quad \Rightarrow \quad \text{ReLU}'(h_{\text{net},2}) = 1 \\
&\quad \delta^{(1)}_2 = (-0.5116) \times (-0.3) \times 1 = 0.15348
\end{align*}

\[
\delta^{(1)} = [0, \; 0.15348]
\]

\separator
\section{گام ۳: انتشار خطا به لایه ورودی (\en{Embedding} سند)}

گرادیان تابع هزینه نسبت به \en{Embedding} سند (\(e_{\text{doc}}\)) به صورت زیر محاسبه می‌شود:

\[
\frac{\partial L}{\partial e_{\text{doc}}} = (\delta^{(1)})^T \cdot W^{(1)}
\]

\[
\frac{\partial L}{\partial e_{\text{doc}}} = 
\begin{bmatrix}
0 & 0.15348
\end{bmatrix}
\cdot
\begin{bmatrix}
0.8 & -0.2 \\
0.1 & 0.6
\end{bmatrix}
\]

\[
\frac{\partial L}{\partial e_{\text{doc},1}} = 0 \times 0.8 + 0.15348 \times 0.1 = 0.015348
\]

\[
\frac{\partial L}{\partial e_{\text{doc},2}} = 0 \times (-0.2) + 0.15348 \times 0.6 = 0.092088
\]

\[
\frac{\partial L}{\partial e_{\text{doc}}} = [0.015348, \; 0.092088]
\]

\separator
\section{گام ۴: انتشار خطا به \en{Embedding} هر کلمه}

از آنجا که \(e_{\text{doc}}\) میانگین \en{Embedding} دو کلمه است:

\[
e_{\text{doc}} = \frac{1}{2}(e_1 + e_2)
\]

بنابراین با استفاده از قانون زنجیره:

\[
\frac{\partial L}{\partial e_1} = \frac{\partial L}{\partial e_{\text{doc}}} \cdot \frac{\partial e_{\text{doc}}}{\partial e_1} = \frac{\partial L}{\partial e_{\text{doc}}} \cdot \frac{1}{2}
\]

\[
\frac{\partial L}{\partial e_2} = \frac{\partial L}{\partial e_{\text{doc}}} \cdot \frac{\partial e_{\text{doc}}}{\partial e_2} = \frac{\partial L}{\partial e_{\text{doc}}} \cdot \frac{1}{2}
\]

\begin{align*}
\frac{\partial L}{\partial e_1} &= \frac{1}{2} \times [0.015348, \; 0.092088] = [0.007674, \; 0.046044] \\[4pt]
\frac{\partial L}{\partial e_2} &= \frac{1}{2} \times [0.015348, \; 0.092088] = [0.007674, \; 0.046044]
\end{align*}

\separator
\section{جمع‌بندی و پاسخ نهایی}

\begin{center}
\fbox{\parbox{0.9\linewidth}{\centering
\textbf{پاسخ نهایی:}
\begin{itemize}
    \item \textbf{گرادیان نسبت به \en{Embedding} کلمه ۱ (ایندکس ۱):}
    \[
    \frac{\partial L}{\partial e_1} = [0.007674, \; 0.046044]
    \]
    \item \textbf{گرادیان نسبت به \en{Embedding} کلمه ۲ (ایندکس ۲):}
    \[
    \frac{\partial L}{\partial e_2} = [0.007674, \; 0.046044]
    \]
\end{itemize}
\textbf{تفسیر:} از آنجا که روش \en{pooling} میانگین گیری بوده، گرادیان به دو کلمه به طور مساوی تقسیم شده است. همچنین به دلیل صفر بودن مشتق \en{ReLU} برای نورون اول لایه پنهان، گرادیان فقط از نورون دوم عبور کرده است.
}}
\end{center}

\thickseparator

\end{document}